% --- Archon physics formalization source begin ---
% archon:physics
% archon:covers PhyXMiniProblems/problem_phyx_mini_0077.lean
% archon:source-report reports/phyx_mini/problem_phyx_mini_0077.source.json

\chapter{Physics problem phyx\_mini\_0077}
\label{ch:PhyXMiniProblems_problem_phyx_mini_0077}

\paragraph{Problem source.}
\#\# Question

What is the height of the final image?

\#\# Answer choices

- A: A: 1.4cm

- B: B: 2.6cm

- C: C: 2.7cm

- D: D: 1.95cm

\#\# Figure caption (auxiliary; use the image as primary evidence)

The image is a diagram illustrating an optical setup with a lens and a mirror.

1. **Objects**:
   - A vertical green arrow on the left represents an object with a height labeled as "1.0 cm".
   - A convex lens in the center with its principal axis aligned horizontally.
   - A concave mirror on the right side.

2. **Text and Labels**:
   - Above the lens, the focal length is indicated as \textbackslash{}( f\_1 = 10 \textbackslash{}, \textbackslash{}text\{cm\} \textbackslash{}).
   - Above the mirror, the focal length is indicated as \textbackslash{}( f\_2 = -30 \textbackslash{}, \textbackslash{}text\{cm\} \textbackslash{}).
   - Below the horizontal line, distances are labeled: "5.0 cm" between the arrow and the lens, and another "5.0 cm" between the lens and the mirror.

3. **Relationships**:
   - The setup aligns the arrow, lens, and mirror along a central horizontal line representing the principal axis.
   - Distances between the components are specified for clarity in positional relationships.

Overall, the diagram is a simple representation of an optical system demonstrating how light moves through the lens and reflects off the mirror.

\#\# Dataset metadata

Geometrical Optics; Spatial Relation Reasoning, Multi-Formula Reasoning

\paragraph{Recorded answer/context.}
C: C: 2.7cm

\paragraph{Figure/image path.}
/root/proposal\_for\_physic/science-mango/phyx\_mini\_run/phyx\_data/test\_image/77.png

\paragraph{Formalization target.}
create a compiling Lean file with sorry bodies at `PhyXMiniProblems/problem\_phyx\_mini\_0077.lean`. The Lean declarations must preserve the physical quantities, dimensions or dimensional roles, figure labels, governing-law hypotheses, and final relation expressed by this problem.
Use Mathlib/Physlib names found through LeanExplore where available. If a physics API is missing, introduce faithful local abstractions rather than scalar placeholder aliases.

\begin{theorem}[Physics formalization target]
\label{thm:physics:phyx_mini_0077:target}
\lean{PhyXMiniProblems.ProblemPhyXMini0077.problem_phyx_mini_0077}
\uses{def:physics:phyx-mini-0077:phyxminiproblems-problemphyxmini0077-lengthincentimeters, def:physics:phyx-mini-0077:phyxminiproblems-problemphyxmini0077-lensmirrorroundtrip, def:physics:phyx-mini-0077:phyxminiproblems-problemphyxmini0077-matchesfigurereadouts, def:physics:phyx-mini-0077:phyxminiproblems-problemphyxmini0077-hasphysicallengthmagnitudes, def:physics:phyx-mini-0077:phyxminiproblems-problemphyxmini0077-satisfiesparaxialroundtripimaginglaw, def:physics:phyx-mini-0077:phyxminiproblems-problemphyxmini0077-matchesdisplayedimageheight}
The assigned autoformalize agent should translate this physics problem into Lean declarations in the covered file, with theorem and lemma proof bodies written as `by sorry`.
\end{theorem}
\begin{proof}
This is an autoformalization task, not a proof task. Produce statements that can later be proved by the physics prover without weakening the physical model.
\end{proof}

% --- Archon declaration topology begin ---
\section{Declaration topology}
\begin{definition}[Length Magnitude]
  \label{def:physics:phyx-mini-0077:phyxminiproblems-problemphyxmini0077-lengthmagnitude}
  \lean{PhyXMiniProblems.ProblemPhyXMini0077.LengthMagnitude}
  A nonnegative physical length magnitude, independent of the chosen units
\end{definition}
\begin{proof}
  This is the declared model definition of Length Magnitude.
\end{proof}
\begin{definition}[length In Centimeters]
  \label{def:physics:phyx-mini-0077:phyxminiproblems-problemphyxmini0077-lengthincentimeters}
  \lean{PhyXMiniProblems.ProblemPhyXMini0077.lengthInCentimeters}
  \uses{def:physics:phyx-mini-0077:phyxminiproblems-problemphyxmini0077-lengthmagnitude}
  The scalar readout of a physical length magnitude in centimeters
\end{definition}
\begin{proof}
  This is the declared model definition of length In Centimeters.
\end{proof}
\begin{definition}[Figure Label]
  \label{def:physics:phyx-mini-0077:phyxminiproblems-problemphyxmini0077-figurelabel}
  \lean{PhyXMiniProblems.ProblemPhyXMini0077.FigureLabel}
  Labels for the three physical elements explicitly visible in the figure
\end{definition}
\begin{proof}
  This is the declared model definition of Figure Label.
\end{proof}
\begin{definition}[Lens Kind]
  \label{def:physics:phyx-mini-0077:phyxminiproblems-problemphyxmini0077-lenskind}
  \lean{PhyXMiniProblems.ProblemPhyXMini0077.LensKind}
  The optical classification of a thin lens
\end{definition}
\begin{proof}
  This is the declared model definition of Lens Kind.
\end{proof}
\begin{definition}[Mirror Kind]
  \label{def:physics:phyx-mini-0077:phyxminiproblems-problemphyxmini0077-mirrorkind}
  \lean{PhyXMiniProblems.ProblemPhyXMini0077.MirrorKind}
  The optical classification of a spherical mirror
\end{definition}
\begin{proof}
  This is the declared model definition of Mirror Kind.
\end{proof}
\begin{definition}[Arrow Orientation]
  \label{def:physics:phyx-mini-0077:phyxminiproblems-problemphyxmini0077-arroworientation}
  \lean{PhyXMiniProblems.ProblemPhyXMini0077.ArrowOrientation}
  Orientation relative to the upright source arrow
\end{definition}
\begin{proof}
  This is the declared model definition of Arrow Orientation.
\end{proof}
\begin{definition}[Optical Approximation]
  \label{def:physics:phyx-mini-0077:phyxminiproblems-problemphyxmini0077-opticalapproximation}
  \lean{PhyXMiniProblems.ProblemPhyXMini0077.OpticalApproximation}
  The approximation used to model the optical system
\end{definition}
\begin{proof}
  This is the declared model definition of Optical Approximation.
\end{proof}
\begin{definition}[Object Arrow]
  \label{def:physics:phyx-mini-0077:phyxminiproblems-problemphyxmini0077-objectarrow}
  \lean{PhyXMiniProblems.ProblemPhyXMini0077.ObjectArrow}
  \uses{def:physics:phyx-mini-0077:phyxminiproblems-problemphyxmini0077-lengthmagnitude, def:physics:phyx-mini-0077:phyxminiproblems-problemphyxmini0077-figurelabel, def:physics:phyx-mini-0077:phyxminiproblems-problemphyxmini0077-arroworientation}
  The source arrow and its physical height
\end{definition}
\begin{proof}
  This is the declared model definition of Object Arrow.
\end{proof}
\begin{definition}[Thin Lens]
  \label{def:physics:phyx-mini-0077:phyxminiproblems-problemphyxmini0077-thinlens}
  \lean{PhyXMiniProblems.ProblemPhyXMini0077.ThinLens}
  \uses{def:physics:phyx-mini-0077:phyxminiproblems-problemphyxmini0077-lengthmagnitude, def:physics:phyx-mini-0077:phyxminiproblems-problemphyxmini0077-figurelabel, def:physics:phyx-mini-0077:phyxminiproblems-problemphyxmini0077-lenskind}
  The thin lens carrying the printed focal label f₁
\end{definition}
\begin{proof}
  This is the declared model definition of Thin Lens.
\end{proof}
\begin{definition}[Spherical Mirror]
  \label{def:physics:phyx-mini-0077:phyxminiproblems-problemphyxmini0077-sphericalmirror}
  \lean{PhyXMiniProblems.ProblemPhyXMini0077.SphericalMirror}
  \uses{def:physics:phyx-mini-0077:phyxminiproblems-problemphyxmini0077-lengthmagnitude, def:physics:phyx-mini-0077:phyxminiproblems-problemphyxmini0077-figurelabel, def:physics:phyx-mini-0077:phyxminiproblems-problemphyxmini0077-mirrorkind}
  The spherical mirror carrying the printed focal label f₂
\end{definition}
\begin{proof}
  This is the declared model definition of Spherical Mirror.
\end{proof}
\begin{definition}[Final Image]
  \label{def:physics:phyx-mini-0077:phyxminiproblems-problemphyxmini0077-finalimage}
  \lean{PhyXMiniProblems.ProblemPhyXMini0077.FinalImage}
  \uses{def:physics:phyx-mini-0077:phyxminiproblems-problemphyxmini0077-lengthmagnitude, def:physics:phyx-mini-0077:phyxminiproblems-problemphyxmini0077-arroworientation}
  The final image is not shown in the setup figure. Its height, orientation, and location on the returning-light side are unknown parts of the model, not definitions made from the recorded answer
\end{definition}
\begin{proof}
  This is the declared model definition of Final Image.
\end{proof}
\begin{definition}[Lens Mirror Round Trip]
  \label{def:physics:phyx-mini-0077:phyxminiproblems-problemphyxmini0077-lensmirrorroundtrip}
  \lean{PhyXMiniProblems.ProblemPhyXMini0077.LensMirrorRoundTrip}
  \uses{def:physics:phyx-mini-0077:phyxminiproblems-problemphyxmini0077-lengthmagnitude, def:physics:phyx-mini-0077:phyxminiproblems-problemphyxmini0077-opticalapproximation, def:physics:phyx-mini-0077:phyxminiproblems-problemphyxmini0077-objectarrow, def:physics:phyx-mini-0077:phyxminiproblems-problemphyxmini0077-thinlens, def:physics:phyx-mini-0077:phyxminiproblems-problemphyxmini0077-sphericalmirror, def:physics:phyx-mini-0077:phyxminiproblems-problemphyxmini0077-finalimage}
  The dimensionful data and labeled components of the pictured round trip
\end{definition}
\begin{proof}
  This is the declared model definition of Lens Mirror Round Trip.
\end{proof}
\begin{definition}[signed Arrow Height In Centimeters]
  \label{def:physics:phyx-mini-0077:phyxminiproblems-problemphyxmini0077-signedarrowheightincentimeters}
  \lean{PhyXMiniProblems.ProblemPhyXMini0077.signedArrowHeightInCentimeters}
  \uses{def:physics:phyx-mini-0077:phyxminiproblems-problemphyxmini0077-lengthmagnitude, def:physics:phyx-mini-0077:phyxminiproblems-problemphyxmini0077-lengthincentimeters, def:physics:phyx-mini-0077:phyxminiproblems-problemphyxmini0077-arroworientation}
  Signed transverse height: upright is positive and inverted is negative
\end{definition}
\begin{proof}
  This is the declared model definition of signed Arrow Height In Centimeters.
\end{proof}
\begin{definition}[signed Lens Focal Length In Centimeters]
  \label{def:physics:phyx-mini-0077:phyxminiproblems-problemphyxmini0077-signedlensfocallengthincentimeters}
  \lean{PhyXMiniProblems.ProblemPhyXMini0077.signedLensFocalLengthInCentimeters}
  \uses{def:physics:phyx-mini-0077:phyxminiproblems-problemphyxmini0077-lengthincentimeters, def:physics:phyx-mini-0077:phyxminiproblems-problemphyxmini0077-thinlens}
  The lens sign is its outgoing focal displacement on the initial left-to-right pass. Thus the converging lens in the figure has positive signed focal length
\end{definition}
\begin{proof}
  This is the declared model definition of signed Lens Focal Length In Centimeters.
\end{proof}
\begin{definition}[signed Mirror Focal Length In Centimeters]
  \label{def:physics:phyx-mini-0077:phyxminiproblems-problemphyxmini0077-signedmirrorfocallengthincentimeters}
  \lean{PhyXMiniProblems.ProblemPhyXMini0077.signedMirrorFocalLengthInCentimeters}
  \uses{def:physics:phyx-mini-0077:phyxminiproblems-problemphyxmini0077-lengthincentimeters, def:physics:phyx-mini-0077:phyxminiproblems-problemphyxmini0077-sphericalmirror}
  The mirror sign is the focal parameter in the unfolded paraxial ray-transfer law. It is positive for a focusing concave mirror and negative for a diverging convex mirror. The primary figure shows the latter and prints f₂ = -30 cm
\end{definition}
\begin{proof}
  This is the declared model definition of signed Mirror Focal Length In Centimeters.
\end{proof}
\begin{definition}[Paraxial Ray]
  \label{def:physics:phyx-mini-0077:phyxminiproblems-problemphyxmini0077-paraxialray}
  \lean{PhyXMiniProblems.ProblemPhyXMini0077.ParaxialRay}
  A scalar paraxial-ray state in a local propagation coordinate
\end{definition}
\begin{proof}
  This is the declared model definition of Paraxial Ray.
\end{proof}
\begin{definition}[propagate Paraxial Ray]
  \label{def:physics:phyx-mini-0077:phyxminiproblems-problemphyxmini0077-propagateparaxialray}
  \lean{PhyXMiniProblems.ProblemPhyXMini0077.propagateParaxialRay}
  \uses{def:physics:phyx-mini-0077:phyxminiproblems-problemphyxmini0077-paraxialray}
  Free paraxial propagation through the given axial distance
\end{definition}
\begin{proof}
  This is the declared model definition of propagate Paraxial Ray.
\end{proof}
\begin{definition}[traverse Thin Lens]
  \label{def:physics:phyx-mini-0077:phyxminiproblems-problemphyxmini0077-traversethinlens}
  \lean{PhyXMiniProblems.ProblemPhyXMini0077.traverseThinLens}
  \uses{def:physics:phyx-mini-0077:phyxminiproblems-problemphyxmini0077-thinlens, def:physics:phyx-mini-0077:phyxminiproblems-problemphyxmini0077-signedlensfocallengthincentimeters, def:physics:phyx-mini-0077:phyxminiproblems-problemphyxmini0077-paraxialray}
  Ideal thin-lens refraction in the local propagation coordinate. This is the ray form of the Gaussian thin-lens law and contains no problem-specific answer
\end{definition}
\begin{proof}
  This is the declared model definition of traverse Thin Lens.
\end{proof}
\begin{definition}[reflect From Spherical Mirror Unfolded]
  \label{def:physics:phyx-mini-0077:phyxminiproblems-problemphyxmini0077-reflectfromsphericalmirrorunfolded}
  \lean{PhyXMiniProblems.ProblemPhyXMini0077.reflectFromSphericalMirrorUnfolded}
  \uses{def:physics:phyx-mini-0077:phyxminiproblems-problemphyxmini0077-sphericalmirror, def:physics:phyx-mini-0077:phyxminiproblems-problemphyxmini0077-signedmirrorfocallengthincentimeters, def:physics:phyx-mini-0077:phyxminiproblems-problemphyxmini0077-paraxialray}
  Ideal spherical-mirror reflection after unfolding the reflected half-space. The mirror focal readout already incorporates the factor of two between the curvature radius and the focal length
\end{definition}
\begin{proof}
  This is the declared model definition of reflect From Spherical Mirror Unfolded.
\end{proof}
\begin{definition}[source Ray]
  \label{def:physics:phyx-mini-0077:phyxminiproblems-problemphyxmini0077-sourceray}
  \lean{PhyXMiniProblems.ProblemPhyXMini0077.sourceRay}
  \uses{def:physics:phyx-mini-0077:phyxminiproblems-problemphyxmini0077-lensmirrorroundtrip, def:physics:phyx-mini-0077:phyxminiproblems-problemphyxmini0077-signedarrowheightincentimeters, def:physics:phyx-mini-0077:phyxminiproblems-problemphyxmini0077-paraxialray}
  A ray leaving the top of the source arrow with arbitrary initial slope
\end{definition}
\begin{proof}
  This is the declared model definition of source Ray.
\end{proof}
\begin{definition}[ray At Final Image Plane]
  \label{def:physics:phyx-mini-0077:phyxminiproblems-problemphyxmini0077-rayatfinalimageplane}
  \lean{PhyXMiniProblems.ProblemPhyXMini0077.rayAtFinalImagePlane}
  \uses{def:physics:phyx-mini-0077:phyxminiproblems-problemphyxmini0077-lengthincentimeters, def:physics:phyx-mini-0077:phyxminiproblems-problemphyxmini0077-lensmirrorroundtrip, def:physics:phyx-mini-0077:phyxminiproblems-problemphyxmini0077-paraxialray, def:physics:phyx-mini-0077:phyxminiproblems-problemphyxmini0077-propagateparaxialray, def:physics:phyx-mini-0077:phyxminiproblems-problemphyxmini0077-traversethinlens, def:physics:phyx-mini-0077:phyxminiproblems-problemphyxmini0077-reflectfromsphericalmirrorunfolded, def:physics:phyx-mini-0077:phyxminiproblems-problemphyxmini0077-sourceray}
  The full sequence is object-to-lens propagation, the first lens traversal, lens-to-mirror propagation, reflection, the unfolded return propagation, the second traversal of the same lens, and propagation to the final plane
\end{definition}
\begin{proof}
  This is the declared model definition of ray At Final Image Plane.
\end{proof}
\begin{definition}[Matches Figure Readouts]
  \label{def:physics:phyx-mini-0077:phyxminiproblems-problemphyxmini0077-matchesfigurereadouts}
  \lean{PhyXMiniProblems.ProblemPhyXMini0077.MatchesFigureReadouts}
  \uses{def:physics:phyx-mini-0077:phyxminiproblems-problemphyxmini0077-lengthincentimeters, def:physics:phyx-mini-0077:phyxminiproblems-problemphyxmini0077-lensmirrorroundtrip, def:physics:phyx-mini-0077:phyxminiproblems-problemphyxmini0077-signedlensfocallengthincentimeters, def:physics:phyx-mini-0077:phyxminiproblems-problemphyxmini0077-signedmirrorfocallengthincentimeters}
  Labels, optical kinds, approximation, and numerical readouts taken from the primary figure. This predicate gives no final-image height, distance, or orientation
\end{definition}
\begin{proof}
  This is the declared model definition of Matches Figure Readouts.
\end{proof}
\begin{definition}[Has Physical Length Magnitudes]
  \label{def:physics:phyx-mini-0077:phyxminiproblems-problemphyxmini0077-hasphysicallengthmagnitudes}
  \lean{PhyXMiniProblems.ProblemPhyXMini0077.HasPhysicalLengthMagnitudes}
  \uses{def:physics:phyx-mini-0077:phyxminiproblems-problemphyxmini0077-lengthincentimeters, def:physics:phyx-mini-0077:phyxminiproblems-problemphyxmini0077-lensmirrorroundtrip}
  Strict positivity selects the nondegenerate physical branch. It does not assign any numerical value to the unknown final-image quantities
\end{definition}
\begin{proof}
  This is the declared model definition of Has Physical Length Magnitudes.
\end{proof}
\begin{definition}[Satisfies Paraxial Round Trip Imaging Law]
  \label{def:physics:phyx-mini-0077:phyxminiproblems-problemphyxmini0077-satisfiesparaxialroundtripimaginglaw}
  \lean{PhyXMiniProblems.ProblemPhyXMini0077.SatisfiesParaxialRoundTripImagingLaw}
  \uses{def:physics:phyx-mini-0077:phyxminiproblems-problemphyxmini0077-lensmirrorroundtrip, def:physics:phyx-mini-0077:phyxminiproblems-problemphyxmini0077-signedarrowheightincentimeters, def:physics:phyx-mini-0077:phyxminiproblems-problemphyxmini0077-rayatfinalimageplane}
  Paraxial image formation means that every ray emitted from the top of the source reaches a common signed transverse point in the final plane. This is a general convergence law; it does not specify the requested numerical height
\end{definition}
\begin{proof}
  This is the declared model definition of Satisfies Paraxial Round Trip Imaging Law.
\end{proof}
\begin{lemma}[final Image Distance In Centimeters eq thirty]
  \label{lem:physics:phyx-mini-0077:phyxminiproblems-problemphyxmini0077-finalimagedistanceincentimeters-eq-thirty}
  \lean{PhyXMiniProblems.ProblemPhyXMini0077.finalImageDistanceInCentimeters_eq_thirty}
  \uses{def:physics:phyx-mini-0077:phyxminiproblems-problemphyxmini0077-lengthincentimeters, def:physics:phyx-mini-0077:phyxminiproblems-problemphyxmini0077-lensmirrorroundtrip, def:physics:phyx-mini-0077:phyxminiproblems-problemphyxmini0077-matchesfigurereadouts, def:physics:phyx-mini-0077:phyxminiproblems-problemphyxmini0077-hasphysicallengthmagnitudes, def:physics:phyx-mini-0077:phyxminiproblems-problemphyxmini0077-satisfiesparaxialroundtripimaginglaw}
  The zero-slope and unit-slope source rays determine the common image plane to be 30 cm from the lens on the returning-light side
\end{lemma}
\begin{proof}
  The conclusion follows from the governing relations and figure readouts named in the statement of final Image Distance In Centimeters eq thirty.
\end{proof}
\begin{lemma}[signed Final Image Height In Centimeters eq]
  \label{lem:physics:phyx-mini-0077:phyxminiproblems-problemphyxmini0077-signedfinalimageheightincentimeters-eq}
  \lean{PhyXMiniProblems.ProblemPhyXMini0077.signedFinalImageHeightInCentimeters_eq}
  \uses{def:physics:phyx-mini-0077:phyxminiproblems-problemphyxmini0077-lensmirrorroundtrip, def:physics:phyx-mini-0077:phyxminiproblems-problemphyxmini0077-signedarrowheightincentimeters, def:physics:phyx-mini-0077:phyxminiproblems-problemphyxmini0077-matchesfigurereadouts, def:physics:phyx-mini-0077:phyxminiproblems-problemphyxmini0077-hasphysicallengthmagnitudes, def:physics:phyx-mini-0077:phyxminiproblems-problemphyxmini0077-satisfiesparaxialroundtripimaginglaw}
  At the common plane the signed transverse height is -8/3 cm; the negative sign records inversion, while the physical height magnitude is 8/3 cm
\end{lemma}
\begin{proof}
  The conclusion follows from the governing relations and figure readouts named in the statement of signed Final Image Height In Centimeters eq.
\end{proof}
\begin{definition}[Answer Choice]
  \label{def:physics:phyx-mini-0077:phyxminiproblems-problemphyxmini0077-answerchoice}
  \lean{PhyXMiniProblems.ProblemPhyXMini0077.AnswerChoice}
  Labels of the four displayed image-height answer choices
\end{definition}
\begin{proof}
  This is the declared model definition of Answer Choice.
\end{proof}
\begin{definition}[displayed Image Height In Centimeters]
  \label{def:physics:phyx-mini-0077:phyxminiproblems-problemphyxmini0077-displayedimageheightincentimeters}
  \lean{PhyXMiniProblems.ProblemPhyXMini0077.displayedImageHeightInCentimeters}
  \uses{def:physics:phyx-mini-0077:phyxminiproblems-problemphyxmini0077-answerchoice}
  The final-image height in centimeters printed beside each answer label
\end{definition}
\begin{proof}
  This is the declared model definition of displayed Image Height In Centimeters.
\end{proof}
\begin{definition}[displayed Height Tolerance]
  \label{def:physics:phyx-mini-0077:phyxminiproblems-problemphyxmini0077-displayedheighttolerance}
  \lean{PhyXMiniProblems.ProblemPhyXMini0077.displayedHeightTolerance}
  \uses{def:physics:phyx-mini-0077:phyxminiproblems-problemphyxmini0077-answerchoice}
  Half a unit in the final printed decimal place of an answer choice
\end{definition}
\begin{proof}
  This is the declared model definition of displayed Height Tolerance.
\end{proof}
\begin{definition}[Matches Displayed Image Height]
  \label{def:physics:phyx-mini-0077:phyxminiproblems-problemphyxmini0077-matchesdisplayedimageheight}
  \lean{PhyXMiniProblems.ProblemPhyXMini0077.MatchesDisplayedImageHeight}
  \uses{def:physics:phyx-mini-0077:phyxminiproblems-problemphyxmini0077-lengthmagnitude, def:physics:phyx-mini-0077:phyxminiproblems-problemphyxmini0077-lengthincentimeters, def:physics:phyx-mini-0077:phyxminiproblems-problemphyxmini0077-answerchoice, def:physics:phyx-mini-0077:phyxminiproblems-problemphyxmini0077-displayedimageheightincentimeters, def:physics:phyx-mini-0077:phyxminiproblems-problemphyxmini0077-displayedheighttolerance}
  A physical height agrees with the displayed rounded answer
\end{definition}
\begin{proof}
  This is the declared model definition of Matches Displayed Image Height.
\end{proof}
% --- Archon declaration topology end ---
% --- Archon physics formalization source end ---
